\documentclass[11pt,a4paper,sans]{moderncv}

% Стиль оформления moderncv: 'classic' отлично структурирует технические блоки
\moderncvstyle{classic}
\moderncvcolor{blue}

% Поля страницы
\usepackage[scale=0.8]{geometry}

% Настройки шрифтов и локализации для LuaLaTeX (согласно методичке)
\usepackage{fontspec}
\defaultfontfeatures{Renderer=Basic, Ligatures=TeX}
\newfontfamily\cyrillicfonttt{CMU Typewriter Text}
\newfontfamily\cyrillicfont{CMU Sans Serif}
\setmainfont{CMU Sans Serif}
\setsansfont{CMU Sans Serif}
\setmonofont{CMU Typewriter Text}

\usepackage{polyglossia}
\setdefaultlanguage{russian}
\setotherlanguages{english}

% Персональные и контактные данные
\name{Леонид}{Тоц}
\title{Разработчик ПО}
\address{Санкт-Петербург}{Россия}
\phone[mobile]{+7~(911)~796-82-29}
\email{leonid005xc@gmail.com}
\homepage{leonidtots.dev}
\social[github]{huksleva}
\extrainfo{Telegram: @pots135}

\begin{document}

\makecvtitle

\section{О себе}
\cvitem{}{Студент направления «Технологии разработки ПО и обработки больших данных» (бакалавриат). Пишу на Python, создаю API, автоматизирую задачи, работаю с реляционными и нереляционными базами данных. Активно применяю FastAPI, Flask, PostgreSQL, MongoDB, Docker. Открыт к позиции Junior-разработчика.}

\section{Образование}
\cventry{2024--н.\,в.}{Бакалавриат, 09.03.01 Информатика и вычислительная техника}{РГПУ им. А. И. Герцена}{Санкт-Петербург}{}{Технологии разработки ПО и обработки больших данных. Перевод с понижением на курс (зима 2024).}
\cventry{2023--2024}{Бакалавриат (2--3 семестр), 01.03.02 Прикладная математика и информатика}{РГПУ им. А. И. Герцена}{Санкт-Петербург}{}{}
\cventry{2023}{Бакалавриат (1 семестр), 09.03.02 Информационные системы и технологии}{Университет ИТМО}{Санкт-Петербург}{}{}

\section{Проекты и достижения}
\cventry{2026}{\textbf{Tramplin -- AI Career Platform}}{Хакатон IF... ELSE 2026 (2-е место)}{}{}{%
Платформа для студентов и работодателей с ролевым доступом, AI-помощью, двухуровневой модерацией и картой возможностей.\\
\textbf{Стек:} Python, FastAPI, PostgreSQL, Docker, JavaScript, Vite.\\
\textbf{Репозиторий:} \url{https://github.com/huksleva/tramplin}}

\vspace{0.2cm}

\cventry{2025}{\textbf{Epidemic Spread Simulation}}{Учебный / исследовательский проект}{}{}{%
Модель клеточного автомата для распространения инфекционных заболеваний (COVID-19, грипп, корь). Реализована визуализация статистики в реальном времени.\\
\textbf{Стек:} Python, NumPy, Matplotlib.\\
\textbf{Репозиторий:} \url{https://github.com/huksleva}}

\section{Технические навыки}
\cvitem{Языки программирования}{Python, JavaScript, TypeScript, C, C++, C\#}
\cvitem{Бэкенд}{FastAPI, Flask, Node.js, SQLAlchemy, Alembic}
\cvitem{Базы данных}{PostgreSQL, MongoDB, Elasticsearch}
\cvitem{Фронтенд}{HTML5, CSS3, React, Next.js, Tailwind CSS, Vite}
\cvitem{Инфраструктура и DevOps}{Docker, Git, Yandex Cloud, Linux}

\end{document}